% GENERATED FILE. Edit canonical lessons, then run npm run generate:books.
% canonical-source-hash: fnv1a64:93f2bfcd9f0a6d99
% canonical-lessons: GU-C205-relve, GU-C205-jagya, GU-C205-utaru, GU-C205-thobhvu, GU-C205-uutarvu

\chapter{A seat on the railway}
\label{ch:gu-a-seat-on-the-railway}
\hlchaptermodality{205}

\begin{chapteropening}
\textbf{By the end of this chapter:} \emph{I can say where I sat, who travelled with me and where we got off.}

Complete the last lesson of chapter 205: I can say where I sat, who travelled with me and where we got off.
\end{chapteropening}

\section[\texorpdfstring{\gu{રેલવે} (relve)}{રેલવે (relve)}]{\gu{રેલવે} (relve) — railway}

\label{lesson:GU-C205-relve}



\begin{quote}
Before the new one: say the Gujarati for to start off, then the Gujarati for to leave.
\end{quote}

\subsection*{You'll want to know: \gu{રેલવે}}
\textbf{\gu{રેલવે}} — \emph{relve} — \textquotedblleft{}railway\textquotedblright{}.

\begin{quote}
\textbf{\gu{રેલવે} \gu{સ્ટેશન} \gu{દૂર} \gu{હતું}.} — \emph{relve sṭeśan dūr hatũ} — \textquotedblleft{}The railway station was far away.\textquotedblright{}
\end{quote}

\begin{cousinweb}
\textbf{\gu{રેલવે}} (\emph{relve}) is a loan from English \emph{railway}.
\end{cousinweb}

\begin{grammarlens}[title={What you've built}]
One more word for everyday life at A2.
\end{grammarlens}

\subsection*{Your turn}
\begin{itemize}
\raggedright
  \item \emph{Say it:} \emph{relve}
  \item \emph{Say it:} \emph{relve}, once more
  \item \emph{Say it:} \emph{relve sṭeśan dūr hatũ} — The railway station was far away.
  \item \emph{Recall:} say the Gujarati for to start off, then the Gujarati for to leave, then say \emph{relve} again
\end{itemize}

\begin{tcolorbox}[breakable,colback=teal!4,colframe=teal!35!black,title={Before you move on}]
What does \gu{રેલવે} mean? (\textquotedblleft{}Railway\textquotedblright{}.) Say it once more.
\end{tcolorbox}
\section[\texorpdfstring{\gu{જગ્યા} (jagyā)}{જગ્યા (jagyā)}]{\gu{જગ્યા} (jagyā) — seat, place, room}

\label{lesson:GU-C205-jagya}



\begin{quote}
Before the new one: say the Gujarati for to leave, then the Gujarati for railway.
\end{quote}

\subsection*{You'll want to know: \gu{જગ્યા}}
\textbf{\gu{જગ્યા}} — \emph{jagyā} — \textquotedblleft{}seat, place, room\textquotedblright{}.

\begin{quote}
\textbf{\gu{ટ્રેનમાં} \gu{જગ્યા} \gu{ન} \gu{હતી}.} — \emph{ṭrenmā̃ jagyā na hatī} — \textquotedblleft{}There was no seat on the train.\textquotedblright{}
\end{quote}

\begin{cousinweb}
\textbf{\gu{જગ્યા}} (\emph{jagyā}) is a loan, usually traced to Persian \emph{jāygāh}, \textquotedblleft{}place\textquotedblright{}. CDIAL has no entry for it.
\end{cousinweb}

\begin{grammarlens}[title={What you've built}]
One more word for everyday life at A2.
\end{grammarlens}

\subsection*{Your turn}
\begin{itemize}
\raggedright
  \item \emph{Say it:} \emph{jagyā}
  \item \emph{Say it:} \emph{jagyā}, once more
  \item \emph{Say it:} \emph{ṭrenmā̃ jagyā na hatī} — There was no seat on the train.
  \item \emph{Recall:} say the Gujarati for to leave, then the Gujarati for railway, then say \emph{jagyā} again
\end{itemize}

\begin{tcolorbox}[breakable,colback=teal!4,colframe=teal!35!black,title={Before you move on}]
What does \gu{જગ્યા} mean? (\textquotedblleft{}Seat, place, room\textquotedblright{}.) Say it once more.
\end{tcolorbox}
\section[\texorpdfstring{\gu{ઉતારુ} (utāru)}{ઉતારુ (utāru)}]{\gu{ઉતારુ} (utāru) — passenger}

\label{lesson:GU-C205-utaru}



\begin{quote}
Before the new one: say the Gujarati for railway, then the Gujarati for seat.
\end{quote}

\subsection*{You'll want to know: \gu{ઉતારુ}}
\textbf{\gu{ઉતારુ}} — \emph{utāru} — \textquotedblleft{}passenger\textquotedblright{}.

\begin{quote}
\textbf{\gu{બસમાં} \gu{ઘણા} \gu{ઉતારુ} \gu{હતા}.} — \emph{basmā̃ ghaṇā utāru hatā} — \textquotedblleft{}There were many passengers on the bus.\textquotedblright{}
\end{quote}

\begin{cousinweb}
\textbf{\gu{ઉતારુ}} (\emph{utāru}) is built on \emph{utār}, which CDIAL glosses \textquotedblleft{}ford, fare for journey\textquotedblright{}, with the ending \emph{-u}. CDIAL lists \emph{utār} under Sanskrit \emph{uttāra}, \textquotedblleft{}landing, delivery\textquotedblright{} (CDIAL 1791).
\end{cousinweb}

\begin{grammarlens}[title={What you've built}]
One more word for everyday life at A2.
\end{grammarlens}

\subsection*{Your turn}
\begin{itemize}
\raggedright
  \item \emph{Say it:} \emph{utāru}
  \item \emph{Say it:} \emph{utāru}, once more
  \item \emph{Say it:} \emph{basmā̃ ghaṇā utāru hatā} — There were many passengers on the bus.
  \item \emph{Recall:} say the Gujarati for railway, then the Gujarati for seat, then say \emph{utāru} again
\end{itemize}

\begin{tcolorbox}[breakable,colback=teal!4,colframe=teal!35!black,title={Before you move on}]
What does \gu{ઉતારુ} mean? (\textquotedblleft{}Passenger\textquotedblright{}.) Say it once more.
\end{tcolorbox}
\section[\texorpdfstring{\gu{થોભવું} (thobhvũ)}{થોભવું (thobhvũ)}]{\gu{થોભવું} (thobhvũ) — to halt, to pause}

\label{lesson:GU-C205-thobhvu}



\begin{quote}
Before the new one: say the Gujarati for seat, then the Gujarati for passenger.
\end{quote}

\subsection*{You'll want to know: \gu{થોભવું}}
\textbf{\gu{થોભવું}} — \emph{thobhvũ} — \textquotedblleft{}to halt, to pause\textquotedblright{}.

\begin{quote}
\textbf{\gu{ટ્રેન} \gu{નાના} \gu{સ્ટેશન} \gu{પર} \gu{થોભી}.} — \emph{ṭren nānā sṭeśan par thobhī} — \textquotedblleft{}The train stopped at a small station.\textquotedblright{}
\end{quote}

\begin{cousinweb}
\textbf{\gu{થોભવું}} (\emph{thobhvũ}) is a tadbhava. It comes from Sanskrit \emph{stōbhatē}, \textquotedblleft{}stops, makes stop\textquotedblright{}, and CDIAL glosses the Gujarati verb \textquotedblleft{}to stop, delay\textquotedblright{} (CDIAL 13724).
\end{cousinweb}

\begin{grammarlens}[title={What you've built}]
One more word for everyday life at A2.
\end{grammarlens}

\subsection*{Your turn}
\begin{itemize}
\raggedright
  \item \emph{Say it:} \emph{thobhvũ}
  \item \emph{Say it:} \emph{thobhvũ}, once more
  \item \emph{Say it:} \emph{ṭren nānā sṭeśan par thobhī} — The train stopped at a small station.
  \item \emph{Recall:} say the Gujarati for seat, then the Gujarati for passenger, then say \emph{thobhvũ} again
\end{itemize}

\begin{tcolorbox}[breakable,colback=teal!4,colframe=teal!35!black,title={Before you move on}]
What does \gu{થોભવું} mean? (\textquotedblleft{}To halt, to pause\textquotedblright{}.) Say it once more.
\end{tcolorbox}
\section[\texorpdfstring{\gu{ઊતરવું} (ūtarvũ)}{ઊતરવું (ūtarvũ)}]{\gu{ઊતરવું} (ūtarvũ) — to get off, to come down}

\label{lesson:GU-C205-uutarvu}



\begin{quote}
Before the new one: say the Gujarati for passenger, then the Gujarati for to halt.
\end{quote}

\subsection*{You'll want to know: \gu{ઊતરવું}}
\textbf{\gu{ઊતરવું}} — \emph{ūtarvũ} — \textquotedblleft{}to get off, to come down\textquotedblright{}.

\begin{quote}
\textbf{\gu{અમે} \gu{બીજા} \gu{સ્ટેશને} \gu{ઊતર્યા}.} — \emph{ame bījā sṭeśane ūtaryā} — \textquotedblleft{}We got off at the second station.\textquotedblright{}
\end{quote}

\begin{cousinweb}
\textbf{\gu{ઊતરવું}} (\emph{ūtarvũ}) is a tadbhava. It comes from Sanskrit \emph{úttarati}, \textquotedblleft{}passes over; alights, descends\textquotedblright{}, and CDIAL glosses the Gujarati verb \textquotedblleft{}to cross, descend, decrease\textquotedblright{} (CDIAL 1770).
\end{cousinweb}

\begin{grammarlens}[title={What you've built}]
One more word for everyday life at A2.
\end{grammarlens}

\subsection*{Your turn}
\begin{itemize}
\raggedright
  \item \emph{Say it:} \emph{ūtarvũ}
  \item \emph{Say it:} \emph{ūtarvũ}, once more
  \item \emph{Say it:} \emph{ame bījā sṭeśane ūtaryā} — We got off at the second station.
  \item \emph{Recall:} say the Gujarati for passenger, then the Gujarati for to halt, then say \emph{ūtarvũ} again
\end{itemize}

\begin{tcolorbox}[breakable,colback=teal!4,colframe=teal!35!black,title={Before you move on}]
What does \gu{ઊતરવું} mean? (\textquotedblleft{}To get off, to come down\textquotedblright{}.) Say it once more.
\end{tcolorbox}
